% SPDX-License-Identifier: LPPL-1.3c
% Copyright (C) 2026 Christophe Jorssen
% Author and Current Maintainer: Christophe Jorssen <christophe.jorssen@gmail.com>
% LPPL maintenance status: maintained
% This file is part of luafemm. See LICENSE and LICENSES.md for its terms.
\section{Complex permeability and laminated-core losses}
\label{sec:losses}

A conducting core can dissipate energy in two different ways: through resolved
currents in the body, and through its magnetic constitutive law. A stack of
insulated sheets adds a third scale: currents circulate within each sheet,
although they cannot travel freely across the insulation. Meshing every sheet
is often unnecessary. This chapter supplies a linear effective law, including
its loss decomposition, and retains the peak $\exp(\mathrm i\omega t)$ convention.

\subsection{A thinner stack, a different frequency response}
Imagine replacing a core made from 0.5 mm sheets by an otherwise identical
stack of 0.2 mm sheets. Give both stacks the same intrinsic material and the
same stacking factor. A material style is the natural place for these choices:
\begin{codeexample}[code only]
\tikzset{femm/materials/core/.style={
  mu r=300,conductivity=2000000,loss angle=10,
  lamination=stacked,lamination fill=.95,lamination thickness=.5}}
\end{codeexample}
The permeability magnitude is 300, the conductivity is $2\,\mathrm{MS/m}$,
and the intrinsic loss angle is $10^\circ$. These are illustrative parameters,
not a fitted electrical-steel grade. |lamination thickness| is in millimetres,
even when the drawing uses another model unit. The stacking factor is the
fraction of magnetic sheet material; the remaining volume is nonmagnetic
insulation with relative permeability one.

First choose a coarse core mesh and a short list of frequencies. Once the
geometry, signs and circuit are satisfactory, refine the core and the air
around it. A sheet thickness is a \emph{material parameter}: it does not set
triangle size. Mesh refinement resolves variations of the averaged field,
while the effective law resolves the one-dimensional variation through a sheet.
The following example uses ordinary PGF loops, named measurements
and |\femmresultplot| to compare both stacks and separate the losses.

\luafemmexamplebegin{harmonic-materials}

% manual example: harmonic-materials/1

\begin{codeexample}[setup code]
\newbox\Calculation
\tikzset{
  femm/materials/core/.style={mu r=300,conductivity=2000000,loss angle=10,
    lamination=stacked,lamination fill=.95,lamination thickness=\Thickness},
  experiment/.style={femm/problem={frequency=\Frequency,depth=20,
    xmin=-8,xmax=8,ymin=-6,ymax=6,mesh size=.45},
    femm/circuit={name=drive,current=2}},
  /femm/region/coil/.style={material=copper,circuit=drive,winding fill=.6,mesh size=.25}}
\newcommand\Geometry{
  \path[femm/region={material=core,mesh size=.2}]
    (-1,-3) rectangle (1,3);
  \path[femm/region={coil,turns=-20}] (-4.7,-2) rectangle (-3.3,2);
  \path[femm/region={coil,turns=20}] (3.3,-2) rectangle (4.7,2);}
\foreach \Case/\Thickness in {thin/0.2,thick/0.5}{
  \foreach \Frequency in {50,100,200,500,1000,2000}{
    \setbox\Calculation=\hbox{
      \begin{tikzpicture}[experiment,
        femm/cache={file=\jobname-\Case-\Frequency,mode=auto}]
        \Geometry
        \femmselect[name=core,material=core]
        \femmmeasure[name=core-\Case-\Frequency,selection=core,parameter=\Frequency]
        \femmmeasure[name=port-\Case-\Frequency,quantity=circuit,circuit=drive]
      \end{tikzpicture}}}}
\def\Frequency{2000}
\end{codeexample}

The following picture uses the preceding definitions.

% manual example: harmonic-materials/2

\begin{codeexample}[width=0pt]
\foreach \Case/\Thickness in {thin/0.2,thick/0.5}{
  \begin{tikzpicture}[experiment,scale=3.3,
    femm/cache={file=\jobname-\Case-\Frequency,mode=auto}]
    \Geometry\pic {femm geometry};
    \pic[femm/lines=16,femm/phase=0,femm/arrow spacing=14mm] {femm field};
    \node[above] at (0,6) {$d=\Thickness\,\mathrm{mm}$, $f=2\,\mathrm{kHz}$};
    \node[below,font=\small] at (0,-6)
      {$P_{\rm core}=\femmvalue[component=magnetic_loss]{core-\Case-2000}$};
  \end{tikzpicture}}
\end{codeexample}

% manual example: harmonic-materials/3

\begin{codeexample}[setup code]
\def\Thin{core-thin-50,core-thin-100,core-thin-200,core-thin-500,core-thin-1000,core-thin-2000}
\def\Thick{core-thick-50,core-thick-100,core-thick-200,core-thick-500,core-thick-1000,core-thick-2000}
\end{codeexample}

The following picture uses the preceding definitions.

% manual example: harmonic-materials/4

\begin{codeexample}[width=0pt]
\begin{tikzpicture}
\begin{loglogaxis}[width=7.5cm,height=6cm,grid=major,
  xlabel={Frequency (Hz)},ylabel={Core loss (W)},
  legend style={font=\scriptsize,at={(.03,.97)},anchor=north west}]
  \femmresultplot[results={\Thin},component=magnetic_loss,
    plot options={blue!70!black,thick,mark=*}]
  \addlegendentry{$d=0.2\,\mathrm{mm}$}
  \femmresultplot[results={\Thick},component=magnetic_loss,
    plot options={red!70!black,thick,mark=square*}]
  \addlegendentry{$d=0.5\,\mathrm{mm}$}
\end{loglogaxis}
\end{tikzpicture}
\end{codeexample}

% manual example: harmonic-materials/5

\begin{codeexample}[width=0pt]
\begin{tikzpicture}
\begin{loglogaxis}[width=7.5cm,height=6cm,grid=major,
  xlabel={Frequency (Hz)},ylabel={Loss components (W)},
  legend style={font=\scriptsize,at={(.03,.97)},anchor=north west}]
  \femmresultplot[results={\Thick},component=hysteresis_loss,
    plot options={blue!70!black,thick,mark=*}]
  \addlegendentry{Intrinsic magnetic loss}
  \femmresultplot[results={\Thick},component=lamination_loss,
    plot options={red!70!black,dashed,thick,mark=square*}]
  \addlegendentry{Microscopic sheet Joule loss}
\end{loglogaxis}
\end{tikzpicture}
\end{codeexample}

\luafemmexampleend{harmonic-materials}

The excitation is a prescribed peak current at every frequency. Field amplitudes
therefore change with frequency and thickness. Do not interpret the loss curves
as a fixed-flux material-characterisation test. The contour levels in the two
pictures are chosen independently. Microscopic sheet currents are not drawn by
|femm eddy currents|: they have been homogenised and are reported as a loss,
not as a fictitious resolved current density in the drawing plane.

\subsection{Material keys and their scope}
All keys below belong to |/tikz/femm/materials| and may appear inside ordinary
material styles. Material styles are evaluated afresh for each captured region.
Optional empty values inherit a supported catalogue property, then the stated
native default. They are not region keys. Declare all materials and regions
before a solve, measurement, field pic or cached solution is requested.

\Needspace{8\baselineskip}
\begin{key}{/tikz/femm/materials/loss angle=\meta{degrees} (initially empty)}
Common intrinsic magnetic loss angle. The native default is zero. For a linear
material, its relative permeability is
$\mu_r^*=|\mu_r|\exp(-\mathrm i\theta)$: the flux density lags the field.
Accepts PGF expressions and finite values $0\leq\theta<90^\circ$.
|mu r|, or the principal |mu x| and |mu y|, supplies the \emph{magnitude},
not the real part. A negative angle would represent an active medium and is
rejected. This parameter does not reconstruct a time-domain hysteresis loop.
\end{key}
\Needspace{8\baselineskip}
\begin{key}{/tikz/femm/materials/loss angle x=\meta{degrees} (initially empty)}
\begin{key}{/tikz/femm/materials/loss angle y=\meta{degrees} (initially empty)}
Principal-axis loss angles. An unspecified angle inherits |loss angle|, whose
native default is zero. Explicit principal angles take precedence over a common
angle. An explicit common angle also replaces catalogue principal defaults;
explicit principal overrides remain effective. |material angle| rotates both
real and imaginary constitutive tensors together in the field plane.
In axisymmetric geometry, the unrotated axes are $r$ and $z$.
\end{key}
\end{key}
\Needspace{8\baselineskip}
\begin{key}{/tikz/femm/materials/lamination=\meta{choice} (initially inherit)}
A choice handler with |inherit|, |none|, |stacked|, |parallel x| and |parallel y|.
|inherit| uses the catalogue interpretation when present, otherwise |none|.
|none| describes a bulk material: its conductivity participates in resolved
eddy-current diffusion unless the region is a stranded winding.

|stacked| describes insulated sheets parallel to the two-component magnetic
field plane, with stacking normal to that plane. Both in-plane principal
permeabilities receive the effective-sheet correction. It suppresses bulk
transport conductivity in the finite-element equation, because the sheet
currents are already included in the effective law. At zero thickness it
means the ideal infinitely fine insulated-sheet limit. It never means a solid
conducting block of thickness zero.

The existing |parallel x| and |parallel y| choices retain their static mixing
interpretation and remain rejected in harmonic problems. |material angle|
rotates principal axes within the sheets; it does not tilt the sheet normal.
In axisymmetry the sheet approximation is local to the meridional field plane;
check that it is appropriate for the physical construction.
\end{key}
\Needspace{7\baselineskip}
\begin{key}{/tikz/femm/materials/lamination thickness=\meta{millimetres} (initially empty)}
Nonnegative thickness of one conducting sheet, excluding insulation. The native
default is zero. Supply a number in millimetres, not a TeX dimension.
A positive value requires |lamination=stacked|. The value is
independent of |mesh size|, picture scale and model length units. A zero value
or zero conductivity removes microscopic sheet-current loss, while intrinsic
magnetic loss may remain. No individual insulation layer is added to the TikZ
geometry or mesh.
\end{key}
\Needspace{7\baselineskip}
\begin{key}{/tikz/femm/materials/lamination fill=\meta{fraction} (initially empty)}
The native default is one; supported catalogue values may supply another value.
Require $0<f\leq1$. For |stacked|, $f$ weights the magnetic sheets and $1-f$
the nonmagnetic insulation. It is unrelated to |winding fill|, which describes
copper occupancy in a stranded winding. A value other than one still requires
a lamination choice. Static |stacked| materials use the zero-frequency arithmetic
mixture along both in-plane axes.
\end{key}

The model is linear. Native stacked sheets require linear intrinsic permeability.
Loss models on captured nonlinear operating points and in the separate axial
$H_z$ diffusion formulation are rejected. A homogenised sheet region cannot
also prescribe an impressed current, belong to a transport circuit, use a
stranded conductor interpretation, or request an isolated-current constraint.
Drive it magnetically from a separate winding or from boundary data.

At frequency zero, loss angles and sheet thickness do not create losses; the
static solve uses the real intrinsic permeability and static fill mixture.
Native cache signatures retain these properties. Static FEM export rejects
new loss metadata instead of silently stripping it. Supported harmonic FEM/ANS interchange preserves intrinsic loss data;
see Section~\ref{sec:interchange}.

\subsection{Which loss should be plotted?}
Harmonic regional measurements expose the following SI components. All are
cycle averages in watts and work with selections, saved |name| values,
|\femmvalue| and |\femmresultplot|. No additional integration key is needed.
\begin{center}
\begin{tabular}{lp{.63\linewidth}}
\toprule
Component & Meaning\\
\midrule
|joule_loss| & Resolved bulk-current loss and stranded-winding copper loss.\\
|eddy_loss| & The existing resolved induced-current or zero-net redistribution loss; it is a part of the electrical loss, not an extra addend.\\
|hysteresis_loss| & Dissipation of the intrinsic complex permeability, integrated through the microscopic sheet field when applicable.\\
|lamination_loss| & Microscopic Joule loss inside insulated conducting sheets.\\
|magnetic_loss| & |hysteresis_loss| plus |lamination_loss|.\\
|total_loss| & |joule_loss| plus |magnetic_loss|.\\
\bottomrule
\end{tabular}
\end{center}
For an ordinary lossless material, the new magnetic components are zero.
For a homogenised sheet with no transport source, |J|, |Jeddy| and |joule_loss|
are zero at the resolved scale, even when |lamination_loss| is positive.
Do not add |eddy_loss| to |total_loss|, and do not infer that a zero resolved
current means zero heating. The intrinsic term is called |hysteresis_loss|
for consistency with magnetic engineering terminology; its constant-angle
linear law is not a microscopic hysteresis model.

The existing energy observable is
$\tfrac14\int\operatorname{Re}(\widehat{\mathbf B}\cdot
\widehat{\mathbf H}^{\,*})\,dV$.
The historical |coenergy| component remains an alias of this average; it is
not a conservative potential from which to differentiate harmonic forces.
In a source-free exterior with zero Dirichlet potential, driven circuit ports
satisfy
\[
 \sum_b\tfrac12\operatorname{Re}(\widehat V_b\widehat I_b^*)
      =P_{\rm total},\qquad
 \sum_b\tfrac12\operatorname{Im}(\widehat V_b\widehat I_b^*)
      =2\omega W.
\]
Boundary excitation or impressed sources outside the measured ports contribute
additional work. These identities check assembly and integration; they do not
replace mesh/domain convergence or material calibration. The energy observable
is not a general thermodynamic energy formula for an arbitrary dispersive medium.
Profiles of |Hx| and |Hy|, including phase, RMS and instantaneous views, use the
full complex law rather than applying a real permeability to each phasor part.

\subsection{Effective law, microscopic losses and numerical checks}
For one principal direction let $\mu^*=\mu_0\mu_r^*$ be the intrinsic absolute
permeability and let $d$ be the sheet thickness in metres. Define
\[
 k=\sqrt{\mathrm i\omega\sigma\mu^*},\quad
 z=kd/2,\quad F(z)=\frac{\tanh z}{z},\quad
 \mu_{r,\mathrm{eff}}^*=1-f+f\mu_r^*F(z).
\]
The passive square-root branch has nonnegative real part. For $z=0$ the
analytic limit is $F(0)=1$. The solver assembles the inverse effective tensor,
including its imaginary stiffness, and the mass term only for resolved bulk
conductivity. Axisymmetric integration retains the usual radial weights.

The intrinsic loss is not obtained by subtracting a second finite-element
solve. The exact sheet field
$H(s)=H_s\cosh(ks)/\cosh(kd/2)$ supplies its mean square. With
$Q=d^{-1}\int_{-d/2}^{d/2}|H(s)/H_s|^2\,ds$, the two physical loss densities are
\[
 p_h=\frac{f\omega\mu_0(-\operatorname{Im}\mu_r^*)}{2}Q|H_s|^2,
 \qquad
 p_e=\frac{f}{2\sigma d}\int_{-d/2}^{d/2}
                    \left|\frac{dH}{ds}\right|^2ds.
\]
Their sum equals
$\tfrac\omega2\operatorname{Im}(\widehat{\mathbf H}\cdot
\widehat{\mathbf B}^{\,*})$ per homogenised volume. The implementation obtains
the sheet term from this power identity, with an explicit passivity check.
It uses a small-argument series and decaying exponentials to avoid cancellation
and overflow in the thin-sheet and deep-skin limits.

The loss regression checks affine complex anisotropy and rotation,
independent numerical integration of the exact microscopic field, extreme
skin-depth limits, phase covariance, coupled planar/axisymmetric power balance,
cache reconstruction and invalid combinations. Independent validation compares independently meshed xfemm solutions
for intrinsic loss, ordinary and deep-skin sheets, axisymmetry, and a core
excited by neighbouring currents. This follows the linear |LamType=0| and
|Theta_hx|/|Theta_hy| conventions in xfemm's harmonic solver; unsupported
in-plane stacking choices are not guessed.

The full catalogue remains available. Supported linear |LamType=0| records
map principal permeability, phase angles, conductivity, sheet thickness and
fill into this law. Catalogue B/H curves retain their nonlinear static meaning:
adding this linear loss model does not turn them into nonlinear harmonic laws.
The supplied catalogue currently has no linear record with nonzero sheet/loss
metadata; the example therefore defines its assumptions explicitly. Effective square and copper-clad wire laws are described in
Section~\ref{sec:wire-keys}; resolved individual strands and arbitrary foil
windings remain outside this homogenised model.
Section~\ref{sec:harmonic-material-models} describes the separate supported
round-wire laws and lossless nonlinear fundamental approximation.

Cache revision 25 fingerprints the intrinsic axes and their orientation, loss angles, fill and
thickness, in addition to frequency. Older snapshots must be recomputed.
Compatible meshes can be reused after a constitutive change, and frozen caches
reconstruct the effective tensors without solving again. Generic cold/warm/shared
frozen tests exercise these properties in plain LuaTeX, LuaLaTeX and ConTeXt MkIV.
